\documentclass[10pt, aspectratio=169]{beamer}
% Preámbulo modular para presentaciones Beamer (Modelación Estadística)
% Diseño y paleta institucional del Tecnológico de Monterrey

\documentclass[aspectratio=169,xcolor={dvipsnames,table}]{beamer}

\usepackage[utf8]{inputenc}
\usepackage[spanish,mexico]{babel}
\usepackage{amsmath,amssymb,amsfonts,mathtools}
\usepackage{graphicx}
\usepackage{listings}
\usepackage{booktabs}
\usepackage{tikz}
\usetikzlibrary{matrix,arrows,decorations.pathmorphing,shapes.geometric,positioning}

% Paleta Institucional Tecnológico de Monterrey
\definecolor{TecAzul}{HTML}{0039A6}       % Azul TEC: color primario institucional
\definecolor{TecAzulOscuro}{HTML}{1A2E51} % Azul marino oscuro
\definecolor{TecRojo}{HTML}{EC2661}       % Rojo de acento (paleta miTec)
\definecolor{TecGris}{HTML}{646464}       % Gris neutro para comentarios y subtítulos
\definecolor{TecFondoBloque}{HTML}{F4F6F9}% Fondo suave para bloques

% Configuración de Tema Beamer
\usetheme{Boadilla}
\usecolortheme[named=TecAzulOscuro]{structure}
\setbeamercolor{alerted text}{fg=TecRojo}
\setbeamercolor{palette primary}{bg=TecAzulOscuro,fg=white}
\setbeamercolor{palette secondary}{bg=TecAzul,fg=white}
\setbeamercolor{palette tertiary}{bg=TecAzulOscuro!80!black,fg=white}
\setbeamercolor{title}{fg=TecAzulOscuro,bg=TecFondoBloque}
\setbeamercolor{frametitle}{fg=TecAzulOscuro,bg=TecFondoBloque}

% Bloques personalizados elegantes
\setbeamercolor{block title}{bg=TecAzulOscuro,fg=white}
\setbeamercolor{block body}{bg=TecFondoBloque,fg=black}
\setbeamercolor{block title alerted}{bg=TecRojo,fg=white}
\setbeamercolor{block body alerted}{bg=TecRojo!8,fg=black}
\setbeamercolor{block title example}{bg=TecAzul,fg=white}
\setbeamercolor{block body example}{bg=TecAzul!8,fg=black}

% Redondear bordes y añadir sombra sutil
\setbeamertemplate{blocks}[rounded][shadow=true]
\setbeamertemplate{navigation symbols}{}

% Pie de página limpio con número de diapositiva
\setbeamertemplate{footline}{
  \leavevmode%
  \hbox{%
  \begin{beamercolorbox}[wd=.7\paperwidth,ht=2.25ex,dp=1ex,left]{author in head/foot}%
    \hspace*{2ex}\insertshortauthor~~(\insertshortinstitute) \hfill \insertshorttitle
  \end{beamercolorbox}%
  \begin{beamercolorbox}[wd=.3\paperwidth,ht=2.25ex,dp=1ex,right]{date in head/foot}%
    \insertshortdate{}\hspace*{2em}
    \insertframenumber{} / \inserttotalframenumber\hspace*{2ex}
  \end{beamercolorbox}}%
  \vskip0pt%
}

% Configuración de Listings de Python para visualización pedagógica de código
\lstset{
    language=Python,
    basicstyle=\ttfamily\footnotesize,
    keywordstyle=\color{TecAzulOscuro}\bfseries,
    stringstyle=\color{TecRojo},
    commentstyle=\color{TecGris}\itshape,
    showstringspaces=false,
    breaklines=true,
    frame=lines,
    rulecolor=\color{TecAzulOscuro!30},
    backgroundcolor=\color{TecFondoBloque},
    aboveskip=0.5em,
    belowskip=0.5em,
    literate={á}{{\'a}}1 {é}{{\'e}}1 {í}{{\'i}}1 {ó}{{\'o}}1 {ú}{{\'u}}1
             {Á}{{\'A}}1 {É}{{\'E}}1 {Í}{{\'I}}1 {Ó}{{\'O}}1 {Ú}{{\'U}}1
             {ñ}{{\~n}}1 {Ñ}{{\~N}}1 {¿}{{?`}}1 {¡}{{!`}}1
}

% Metadatos institucionales por defecto
\title[Modelación Estadística]{Modelación Estadística para Ciencia de Datos e Ingeniería}
\author[J. Castillo Colmenares]{Juliho David Castillo Colmenares}
\institute[Tec de Monterrey]{Tecnológico de Monterrey \\ Escuela de Ingeniería y Ciencias}
\date{\today}

% Comandos y macros estadísticas compartidas para las presentaciones Beamer (Metropolis)

% Conjuntos numéricos básicos
\newcommand{\R}{\mathbb{R}}
\newcommand{\N}{\mathbb{N}}
\newcommand{\Q}{\mathbb{Q}}
\newcommand{\Z}{\mathbb{Z}}
\newcommand{\set}[1]{\left\{ #1 \right\}}
\newcommand{\abs}[1]{\left|#1\right|}
\newcommand{\norm}[1]{\left\| #1 \right\|}

% Comandos estadísticos y probabilísticos del curso
\newcommand{\Var}{\operatorname{Var}}
\newcommand{\cov}{\operatorname{Cov}}
\newcommand{\corr}{\rho}
\newcommand{\comb}[2]{\begin{pmatrix} #1 \\ #2 \end{pmatrix}}
\newcommand{\s}{\sigma}
\newcommand{\particion}{\mathfrak{P}}
\newcommand{\card}[1]{\#\left( #1 \right)}
\newcommand{\seq}[3]{\set{#1}_{#2}^{#3}}
\newcommand{\E}{\mathbb{E}}
\newcommand{\Prob}{\mathbb{P}}

% Entornos pedagógicos y cajas en español académico natural
\newenvironment{discusion}{%
  \begin{alertblock}{Pregunta de Discusión}%
}{%
  \end{alertblock}%
}

\newenvironment{reflexion}{%
  \begin{alertblock}{Reflexión para el Aula}%
}{%
  \end{alertblock}%
}

\newenvironment{paradoja}{%
  \begin{alertblock}{Paradoja de Exploración}%
}{%
  \end{alertblock}%
}

\newenvironment{motivacion}{%
  \begin{exampleblock}{Motivación: Ciencia de Datos en la Práctica}%
}{%
  \end{exampleblock}%
}

\newenvironment{retoclase}{%
  \begin{block}{Reto Rápido en Equipo}%
}{%
  \end{block}%
}


\title[Distribuciones en Ciencia de Datos]{Sección 03.10: Distribuciones Discretas en Ciencia de Datos}
\subtitle{Ajuste por MLE, Detección de Sobredispersión y Criterios de Información}
\author[J. Castillo Colmenares]{Juliho Castillo Colmenares}
\institute{Tecnológico de Monterrey}
\date{\vspace{-1.2cm}}

\begin{document}

% Slide 1: Portada
\begin{frame}[plain]
  \vspace{-0.8cm}
  \titlepage
\end{frame}

% Slide 2: Hoja de Ruta
\begin{frame}{Hoja de Ruta --- Capítulo 03: Variables Aleatorias Discretas}
  \footnotesize
  ¿Dónde nos encontramos en el desarrollo modular de la teoría? \pause
  \begin{itemize}\itemsep=0.08em
    \item \textbf{03.01} PMF y Soporte \pause
    \item \textbf{03.02} CDF Discreta \pause
    \item \textbf{03.03} Esperanza y Varianza \pause
    \item \textbf{03.04} Bernoulli y Binomial \pause
    \item \textbf{03.05} Geométrica y Binomial Negativa \pause
    \item \textbf{03.06} Hipergeométrica \pause
    \item \textbf{03.07} Poisson \pause
    \item \textbf{03.08} Multinomial \pause
    \item \textbf{03.09} Normal y Aproximación Continua \pause
    \item \color{TecRojo}\textbf{03.10 Distribuciones Discretas en Ciencia de Datos \emph{(Hoy y Cierre del Cap. 03)}}\color{black}
  \end{itemize}
  \vspace{-0.05cm}
  \begin{block}{Objetivo de la Sesión}
    Integrar todas las distribuciones discretas en aplicaciones modernas: ajuste por máxima verosimilitud, detección de sobredispersión, criterios de información para selección de modelos, e inferencia basada en razón de verosimilitud.
  \end{block}
\end{frame}

% Slide 3: Motivación
\begin{frame}{De la Teoría a la Práctica: Distribuciones en Datos Reales}
  \footnotesize
  \begin{motivacion}[¿Por qué ajustar distribuciones a datos empíricos?]
    En la práctica, rara vez sabemos con certeza qué distribución generó nuestros datos. Lo que observamos son conteos (clics, defectos, cancelaciones, llamadas) y necesitamos inferir el modelo que mejor los describe. La \emph{máxima verosimilitud} (MLE) es el método estándar: dada una familia paramétrica, elegimos los parámetros que maximizan la probabilidad de observar los datos.
  \end{motivacion}

  \vspace{0.15cm}
  \pause
  \begin{itemize}\itemsep=0.1cm
    \item \textbf{Ajuste paramétrico:} La Poisson asume equidispersión; si la varianza empírica es mayor, la Binomial Negativa es preferible. \pause
    \item \textbf{Selección de modelos:} AIC, BIC y likelihood ratio test balancean bondad de ajuste y parsimonia. \pause
    \item \textbf{Incertidumbre:} Los intervalos de confianza y el bootstrap cuantifican la variabilidad de los estimadores puntuales.
  \end{itemize}
\end{frame}

% Slide 4: MLE para Poisson
\begin{frame}{Estimación por Máxima Verosimilitud (MLE) para Poisson}
  \footnotesize
  Para datos de conteo $x_{1}, \dots, x_{n}$ asumiendo $X_{i} \sim \text{Pois}(\lambda)$ independientes, la función de verosimilitud es: \pause

  \vspace{0.1cm}
  \begin{block}{Función de Verosimilitud Poisson}
    $$L(\lambda) = \prod_{i=1}^{n} \dfrac{e^{-\lambda}\lambda^{x_{i}}}{x_{i}!} = \dfrac{e^{-n\lambda}\lambda^{\sum x_{i}}}{\prod x_{i}!}$$
    La log-verosimilitud es $\ell(\lambda) = -n\lambda + (\sum x_{i})\log \lambda - \sum \log x_{i}!$.
  \end{block}

  \vspace{0.1cm}
  \pause
  \begin{alertblock}{Derivación del MLE}
    \begin{itemize}\itemsep=0.05cm
      \item Derivando: $\ell'(\lambda) = -n + \dfrac{\sum x_{i}}{\lambda} = 0 \implies \hat{\lambda} = \bar{x}$.
      \item El MLE de Poisson es simplemente la media muestral: la estimación más natural para un parámetro de tasa.
      \item Por el TCL, $\hat{\lambda} \pm z_{0.025}\sqrt{\hat{\lambda}/n}$ es un intervalo de confianza al $95\%$ asintótico.
    \end{itemize}
  \end{alertblock}
\end{frame}

% Slide 5: Detección de Sobredispersión
\begin{frame}{Detección de Sobredispersión: Índice de Dispersión $D = s^{2}/\bar{x}$}
  \footnotesize
  La \emph{sobredispersión} ocurre cuando la varianza muestral excede la media muestral, violando el supuesto de equidispersión de Poisson. El diagnóstico estándar es el \emph{índice de dispersión}: \pause

  \vspace{0.1cm}
  \begin{block}{Estadístico de Diagnóstico}
    $$D = \dfrac{s^{2}}{\bar{x}}, \quad \text{con } s^{2} = \dfrac{1}{n-1}\sum_{i=1}^{n}(x_{i} - \bar{x})^{2}$$
    Bajo $H_{0}$: Poisson (equidispersión), $D \approx 1$. La regla práctica:
  \end{block}

  \vspace{0.1cm}
  \pause
  \begin{alertblock}{Interpretación del Cociente de Dispersión}
    \begin{itemize}\itemsep=0.05cm
      \item $D \approx 1$: equidispersión; Poisson es adecuada.
      \item $D > 1.5$: sobredispersión fuerte; se recomienda Binomial Negativa o regresión de Poisson con varianza robusta.
      \item $D < 0.8$: sub-dispersión; considerar distribución de Conway-Maxwell-Poisson.
    \end{itemize}
  \end{alertblock}
\end{frame}

% Slide 6: Binomial Negativa como Alternativa

% Slide 7: Criterios de Información
\begin{frame}{Criterios de Información: AIC y BIC para Selección de Modelos}
  \footnotesize
  Para comparar modelos con diferente número de parámetros, los criterios de información penalizan la complejidad: \pause

  \vspace{0.1cm}
  \begin{block}{Criterio de Información de Akaike (AIC) y Bayesiano (BIC)}
    $$\text{AIC} = 2k - 2\ell(\hat{\theta}), \quad \text{BIC} = k \log n - 2\ell(\hat{\theta})$$
    donde $k$ es el número de parámetros, $n$ el tamaño muestral, y $\ell(\hat{\theta})$ la log-verosimilitud maximizada.
  \end{block}

  \vspace{0.1cm}
  \pause
  \begin{alertblock}{Regla de Decisión}
    \begin{itemize}\itemsep=0.05cm
      \item Se prefiere el modelo con menor AIC (o BIC).
      \item Diferencias $|\Delta\text{AIC}| < 2$ son insignificantes; $4 < \Delta\text{AIC} < 7$ son sustanciales; $\Delta\text{AIC} > 10$ son muy fuertes.
      \item BIC penaliza más la complejidad que AIC; para muestras grandes favorece modelos más parsimoniosos.
    \end{itemize}
  \end{alertblock}
\end{frame}

% Slide 8: Test de Razón de Verosimilitud
\begin{frame}{Test de Razón de Verosimilitud (Wilks) y Sobredispersión}
  \footnotesize
  El test de razón de verosimilitud compara un modelo restringido $L_{0}$ contra uno completo $L_{1}$: \pause

  \vspace{0.1cm}
  \begin{block}{Estadístico de Razón de Verosimilitud}
    $$\Lambda = -2(\log L_{0}(\hat{\theta}_{0}) - \log L_{1}(\hat{\theta}_{1})) \xrightarrow{d} \chi^{2}_{k_{1} - k_{0}}$$
    Bajo $H_{0}$ (modelo restringido correcto), $\Lambda$ converge en distribución a chi-cuadrado con $k_{1} - k_{0}$ grados de libertad.
  \end{block}

  \vspace{0.1cm}
  \pause
  \begin{alertblock}{Aplicación: Test de Sobredispersión}
    Para $H_{0}$: Poisson vs. $H_{1}$: Binomial Negativa, calculamos:
    $$\Lambda = -2(\log L_{\text{Poisson}}(\hat{\lambda}) - \log L_{\text{BN}}(\hat{r}, \hat{p})) \sim \chi^{2}_{1}$$
    Rechazamos $H_{0}$ si $\Lambda > \chi^{2}_{1, 0.05} = 3.841$, concluyendo sobredispersión significativa.
  \end{alertblock}
\end{frame}

% Slide 9: Intervalos de Confianza
\begin{frame}{Intervalos de Confianza: Wald, Exacto y Bootstrap}
  \footnotesize
  Para cuantificar la incertidumbre en la estimación de $\lambda$, existen tres métodos principales: \pause

  \vspace{0.1cm}
  \begin{columns}[T]
    \begin{column}{0.48\textwidth}
      \begin{block}{Método de Wald}
        \scriptsize
        $\hat{\lambda} \pm z_{0.025}\sqrt{\hat{\lambda}/n}$. Válido para $\hat{\lambda} \ge 30$.
      \end{block}
    \end{column}
    \begin{column}{0.48\textwidth}
      \begin{block}{Método Exacto (Garwood)}
        \scriptsize
        $\left[\dfrac{1}{2n}\chi^{2}_{0.025, 2k}, \dfrac{1}{2n}\chi^{2}_{0.975, 2(k+1)}\right]$ con $k = \sum x_{i}$. Válido para todo $\lambda$.
      \end{block}
    \end{column}
  \end{columns}

  \vspace{0.15cm}
  \pause
  \begin{alertblock}{Bootstrap No Paramétrico}
    \begin{itemize}\itemsep=0.04cm
      \item Muestreo con reemplazo de los datos observados: $B$ muestras, cada una de tamaño $n$.
      \item Estimación del error estándar como la desviación estándar de los B MLEs bootstrap.
      \item Intervalo de confianza por percentiles: $[\hat{\lambda}^{*}_{0.025}, \hat{\lambda}^{*}_{0.975}]$.
    \end{itemize}
  \end{alertblock}
\end{frame}

% Slide 10: Aplicaciones
\begin{frame}{Aplicaciones en Ciencia de Datos y Negocios}
  \footnotesize
  El ajuste de distribuciones discretas tiene aplicaciones directas en analytics empresarial: \pause

  \vspace{0.15cm}
  \begin{itemize}\itemsep=0.12cm
    \item \textbf{Análisis de Churn:} Modelado de cancelaciones de clientes con Poisson y Binomial Negativa para identificar factores de riesgo en cohortes. \pause
    \item \textbf{Detección de Fraude:} Conteo de transacciones sospechosas; la sobredispersión señala comportamiento anormal. \pause
    \item \textbf{Experimentación A/B:} Comparación de tasas de conversión (Binomial) y conteos de eventos (Poisson/NegBin) entre variantes. \pause
    \item \textbf{Confiabilidad Industrial:} Tiempo hasta la falla (Exponencial/Weibull) y número de defectos por unidad (Poisson/NegBin).
  \end{itemize}
\end{frame}

% Slide 12A-1: Lab Python (Bloque 1)
\begin{frame}[fragile]{Laboratorio en Python: MLE y Sobredispersión}
  \scriptsize
  Ajuste de modelos Poisson y Binomial Negativa, comparación vía AIC/BIC (\texttt{03.10\_discrete\_distributions\_data\_science.py}):
  \vspace{0.02cm}
  {\lstset{basicstyle=\fontsize{4.5pt}{5.5pt}\selectfont\ttfamily}
  \lstinputlisting[firstline=15, lastline=40]{../../code/03_variables_aleatorias_discretas/03.10_discrete_distributions_data_science.py}}
\end{frame}

% Slide 12A-2: Lab Python (Bloque 2)
\begin{frame}[fragile]{Laboratorio en Python: Test de Wilks y Sobredispersión}
  \scriptsize
  Aplicación del test de razón de verosimilitud para detectar sobredispersión significativa:
  \vspace{0.02cm}
  {\lstset{basicstyle=\fontsize{4.5pt}{5.5pt}\selectfont\ttfamily}
  \lstinputlisting[firstline=48, lastline=72]{../../code/03_variables_aleatorias_discretas/03.10_discrete_distributions_data_science.py}}
\end{frame}

% Slide 12B: Lab Python (Bloque 3)
\begin{frame}[fragile]{Laboratorio en Python: IC y Bootstrap}
  \scriptsize
  Tres métodos de intervalo de confianza (Wald, exacto, bootstrap) y bootstrap para parámetros de Binomial Negativa:
  \vspace{0.02cm}
  {\lstset{basicstyle=\fontsize{4.5pt}{5.5pt}\selectfont\ttfamily}
  \lstinputlisting[firstline=81, lastline=108]{../../code/03_variables_aleatorias_discretas/03.10_discrete_distributions_data_science.py}}
\end{frame}

% Slide 12C: Salida Terminal
\begin{frame}[fragile]{Laboratorio en Python: Salida en Terminal}
  \scriptsize
  Salida estándar generada al ejecutar el script del laboratorio en Python:
  \vspace{0.02cm}
  \begin{block}{Salida Estándar en Consola}
    \fontsize{4pt}{5pt}\selectfont
    \begin{verbatim}
=== Block 1: MLE Fitting & Overdispersion ===
NegBin(r=2, p=0.4) | mean=2.95 | var=5.24 | D=1.78
  -> Strong overdispersion (D > 1.5): NegBin recommended
Poisson: logLik=-220.16, AIC=442.31, BIC=444.92
NegBin:  logLik=-211.05, AIC=426.10, BIC=431.31
  Delta AIC (P-NB): 16.21 | Delta BIC: 13.60

=== Block 2: Likelihood Ratio Test & Wilks ===
NegBin(r=3, p=0.5) n=200: mean=2.83, var=5.29
  Poisson logLik=-444.31 | NegBin logLik=-421.64
  Lambda = 45.34 | p-value = 0.000000
  Decision: Reject H0. NegBin significantly better.
LR vs overdispersion (n=100):
  r=10.0 (v/m=1.33): LR=4.24, p=0.040
  r=2.0  (v/m=1.58): LR=12.81, p=0.000
  r=1.0  (v/m=2.85): LR=78.66, p=0.000
  r=0.5  (v/m=7.99): LR=310.80, p=0.000

=== Block 3: Confidence Intervals & Bootstrap ===
12 months claims: sample mean=46.50
  Wald 95% CI: [42.64, 50.36]
  Exact 95% CI: [42.72, 50.52]
  Bootstrap SE: 1.97 | Wald SE: 1.97
  Bootstrap 95% CI: [42.75, 50.42]
NegBin bootstrap (n=200 churn data):
  Method of moments: r=2.03, p=0.40
  Bootstrap mean: r=2.13 +- 0.42 | p=0.41 +- 0.05
    \end{verbatim}
  \end{block}
\end{frame}

% Slide 17: Cierre
\begin{frame}{Cierre del Capítulo 03 y Síntesis Final}
  \begin{reflexion}[El Viaje desde Binomial hasta Ciencia de Datos]
    A lo largo de 10 secciones, hemos dominado el ecosistema completo de distribuciones discretas: Binomial, Geométrica, Hipergeométrica, Poisson, Multinomial, y su aproximación Normal. Ahora cerramos con las herramientas de inferencia que conectan estos modelos con datos empíricos: MLE, criterios de información, tests de razón de verosimilitud, e inferencia por bootstrap.
  \end{reflexion}

  \vspace{0.2cm}
  \pause
  \begin{block}{Lecciones del Capítulo 03: Variables Aleatorias Discretas}
    \begin{itemize}\itemsep=0.06cm
      \item \textbf{Familia completa:} Binomial, Geométrica, Hipergeométrica, Poisson, Multinomial, Normal.
      \item \textbf{Equidispersión vs. Sobredispersión:} La Poisson asume $\E[X] = \Var(X)$; la Binomial Negativa relaja este supuesto.
      \item \textbf{Inferencia rigurosa:} AIC/BIC balancean bondad de ajuste y parsimonia; Wilks valida modelos anidados.
    \end{itemize}
  \end{block}
\end{frame}

% Slide 18: Perspectiva Modular
\begin{frame}{Perspectiva Modular --- Hacia el Capítulo 04}
  \scriptsize
  \begin{retoclase}[Próximo Capítulo: Variables Aleatorias Continuas]
    Las distribuciones discretas modelan conteos, pero muchos fenómenos en ingeniería (tiempos de falla, errores de medición, señales) varían continuamente. En el Capítulo 04 introduciremos las \emph{distribuciones continuas fundamentales}: Uniforme, Exponencial, Normal, Gamma, Beta y Weibull, junto con sus aplicaciones en procesos continuos sin memoria y análisis de confiabilidad.
  \end{retoclase}

  \vspace{0.08cm}
  \pause
  \begin{alertblock}{Preparación Recomendada}
    \begin{itemize}\itemsep=0.06cm
      \item Repasa el concepto de función de densidad de probabilidad (PDF) y la diferencia con la PMF.
      \item Reflexiona sobre cómo el tiempo entre eventos Poisson se distribuye exponencialmente (proceso de Poisson).
    \end{itemize}
  \end{alertblock}
\end{frame}

\end{document}
